% Springer Nature sn-jnl layout. Like every other template in this project it adds no content: the sections live in
% sections/, the numbers in latex/, and this file only supplies the class, the author block and the bibliography
% style. Switching journal means writing one of these and nothing else.
%
% The class and its styles are vendored in templates/springer/; TEXINPUTS must include that directory, which
% .latexmkrc sets so that every caller finds it. The class option chooses the bibliography style:
% sn-vancouver-num for a numbered journal, sn-basic or sn-vancouver-ay for author-year.
%
% Springer asks for a single .tex at submission. Flatten this file then, for instance with
%   latexpand paper_springer.tex > submission.tex
% and keep working with the \input version, which is what the pipeline verifies.
\documentclass[pdflatex,sn-vancouver-num]{sn-jnl}

\usepackage{graphicx}
\usepackage{multirow}
\usepackage{amsmath,amssymb,amsfonts}
\usepackage{array}
\usepackage{tabularx}
\usepackage{booktabs}
\usepackage{longtable}
\usepackage{xcolor}

\raggedbottom

% Read in the preamble, not next to the generated macros below: the class expands the running
% title straight away, so these definitions have to exist before \title.
% The front matter as data: title, authors, their e-mail, ORCID and contributions, and the one affiliation.
% No journal's syntax appears here, because every class asks for the same facts in a different shape: one wants
% a single name per author with the address as key=value pairs, another wants the given and family names split
% into separate commands and the address broken into fields of its own, and the reading copy wants one
% \and-separated line. So this file states each fact once and a template renders it, exactly as the abstract and
% the keywords already work. Supporting one more journal means writing the renderers below in the new
% template, and changing nothing here.
%
% A template must define, before it uses \projauthorlist:
%   \projauthor{initials}{given}{family}{email}{orcid}{CRediT roles}      one non-corresponding author
%   \projauthorcorr{...}                                                 same six, for the corresponding author
% ORCID may be empty, so a renderer that prints it guards the field with \projifhasorcid.
%
% THIS IS THE FIRST FILE TO EDIT IN A NEW PAPER. Nothing else holds any of these facts.

\newcommand{\projtitle}{The Title of the Paper, Stated in Full and Without a Subtitle Unless One Earns Its Place}
\newcommand{\projshorttitle}{Short Running Title}
\newcommand{\projrunningauthors}{Family-Name et al., 2026}

\newcommand{\projaffildiv}{Department or Institute}
\newcommand{\projaffilorg}{University}
\newcommand{\projaffilstreet}{Street and Number}
\newcommand{\projaffilcity}{City}
\newcommand{\projaffilstate}{Region}
\newcommand{\projaffilpostcode}{00000}
\newcommand{\projaffilcountry}{Country}

% Not every author has an ORCID on file, and a renderer that prints the field unconditionally would emit an empty
% link for the ones who do not. This runs its second argument only when the first is a non-empty ORCID.
\newcommand{\projifhasorcid}[2]{\def\projorcidarg{#1}\ifx\projorcidarg\empty\else#2\fi}

% The authors in the order they appear on the paper. The corresponding author is marked by which macro the line
% calls, not by a flag, so no layout has to decide what "corresponding" looks like in its own notation.
% The roles are the CRediT taxonomy's own terms: they go into the author block of a class that has one, and into
% the contributions declaration, which is rendered from this same list.
\newcommand{\projauthorlist}{%
	\projauthorcorr{A.B.}{Ada}{Firstauthor}{ada.firstauthor@example.org}{0000-0000-0000-0000}%
		{Conceptualization, Methodology, Software, Formal analysis, Writing -- original draft}%
	\projauthor{C.D.}{Carl}{Secondauthor}{carl.secondauthor@example.org}{}%
		{Validation, Writing -- review \& editing, Supervision}%
}


\begin{document}

\title[\projshorttitle]{\projtitle}

% The front matter in sn-jnl notation: the given and family names split, the corresponding author marked by
% \author*, and the address decomposed into the class's own fields. The facts are in
% sections/manuscript_metadata.tex and appear nowhere else.
\newcommand{\projauthor}[6]{\author[1]{\fnm{#2} \sur{#3}}\email{#4}}
\newcommand{\projauthorcorr}[6]{\author*[1]{\fnm{#2} \sur{#3}}\email{#4}}
\projauthorlist

\affil[1]{\orgdiv{\projaffildiv},
	\orgname{\projaffilorg},
	\orgaddress{\street{\projaffilstreet}, \city{\projaffilcity}, \postcode{\projaffilpostcode},
		\state{\projaffilstate}, \country{\projaffilcountry}}}

% A generated table, typeset at the full text width. Every table in latex/ is a tabularx or a tabular* that
% already fills \textwidth, so no template has to box or scale one; this hook stays so that a journal needing
% something of its own can redefine it in one place.
\newcommand{\projtable}[1]{\input{#1}}

% One declaration: the statement is shared, the heading belongs to the layout. The declarations are
% listed in sections/manuscript_declarations.tex and this is the only place a layout decides how they
% are presented.
\newcommand{\projdeclaration}[2]{\subsection*{#1}\input{#2}}

% Every macro file the pipeline writes, read once so that no document can be missing one. A new step that emits
% numbers adds its file here and nowhere else; tests/test_manuscript_layout.py fails if latex/ holds a
% *_variables.tex that this list does not name.
% Written by scripts/99_demo_variables.py. Do not edit: the next run overwrites it.
\newcommand{\DemoHeadline}{0.74}
\newcommand{\DemoNSamples}{859}



% sn-jnl stores the abstract and the keywords and typesets them with \maketitle, so both must be set before it.
\abstract{% The abstract text only, with no environment of its own. Some classes want it wrapped in an environment and
% others in a command, and a class that typesets the environment where it stands would put the abstract before
% the title, so the choice belongs to the template.
%
% Write it last. One useful order: the problem, what was done, three results, what the reader may not conclude.
% Every number here is a macro from latex/, never a typed digit: \DemoHeadline is an example of one.
The problem this paper addresses, in one sentence that a reader outside the field can follow. What was done, and
on what data. The main result, which came out at \DemoHeadline\ on \DemoNSamples\ samples. The second result. The
third. What the reader should not conclude from any of it.
}
\keywords{% The keyword list only, comma separated, for whatever the template's keyword command or environment is.
first keyword, second keyword, third keyword, fourth keyword, fifth keyword
}

\maketitle

% The manuscript itself, independent of any journal template: the sections in order, each one free of preamble.
% A template supplies the class, the author block and the bibliography style; this file supplies the content.
% The declarations (funding, competing interests and the rest) are not here: a layout wraps them in its own
% way, so it reads sections/manuscript_declarations.tex separately, after this file.
%
% This list is the order of the paper. tests/test_manuscript_layout.py checks that the numbered files are listed
% in ascending order, so renumber the files when you reorder them.
% The Introduction, with no preamble of its own: a template supplies the class, the front matter and the bibliography
% style. Cite with \citep{} and \citet{}; refer to a generated table with \projtable, never \input.
\section{Introduction}

Replace this paragraph. Every number in the text comes from a macro that a script in scripts/ wrote into latex/;
the check step fails the build if the text names a macro that nothing produced. Cite as
\citet{example2026} or \citep{example2026}, from references.bib.

% The Methods, with no preamble of its own: a template supplies the class, the front matter and the bibliography
% style. Cite with \citep{} and \citet{}; refer to a generated table with \projtable, never \input.
\section{Methods}

Replace this paragraph. Every number in the text comes from a macro that a script in scripts/ wrote into latex/;
the check step fails the build if the text names a macro that nothing produced.

% The Results, with no preamble of its own: a template supplies the class, the front matter and the bibliography
% style. Cite with \citep{} and \citet{}; refer to a generated table with \projtable, never \input.
\section{Results}

Replace this paragraph. Every number in the text comes from a macro that a script in scripts/ wrote into latex/;
the check step fails the build if the text names a macro that nothing produced.

\begin{table}[htbp]
	\centering
	\caption{An example of a generated table. The caption is written by hand here, next to the item it describes,
		while the body comes from latex/ and is overwritten on every run; a caption that states a number states it
		with a macro, as this one does: \DemoNSamples{} samples.}
	\label{tab:demo}
	\projtable{latex/demo_table.tex}
\end{table}

% The Discussion, with no preamble of its own: a template supplies the class, the front matter and the bibliography
% style. Cite with \citep{} and \citet{}; refer to a generated table with \projtable, never \input.
\section{Discussion}

Replace this paragraph. Every number in the text comes from a macro that a script in scripts/ wrote into latex/;
the check step fails the build if the text names a macro that nothing produced.

% The Conclusions, with no preamble of its own: a template supplies the class, the front matter and the bibliography
% style. Cite with \citep{} and \citet{}; refer to a generated table with \projtable, never \input.
\section{Conclusions}

Replace this paragraph. Every number in the text comes from a macro that a script in scripts/ wrote into latex/;
the check step fails the build if the text names a macro that nothing produced.



\backmatter
\section*{Declarations}
% The declarations publishers ask for, in the order they are most often requested. Each one is a file of its own so
% that a statement is written once and every layout reports the same thing; the heading, and whether the set is
% gathered under one of its own, is the layout's business, which is what \projdeclaration supplies. Adding a
% journal means defining that command and reordering these lines, not rewriting the statements.
\projdeclaration{Availability of data and materials}{sections/06_1_Data_Available.tex}
\projdeclaration{Competing interests}{sections/06_2_Competing_Interests.tex}
\projdeclaration{Funding}{sections/06_3_Funding.tex}
\projdeclaration{Authors' contributions}{sections/06_4_Authors_Contributions.tex}
\projdeclaration{Acknowledgements}{sections/06_5_Acknowledgements.tex}
\projdeclaration{Declaration of generative AI use}{sections/06_6_Generative_AI.tex}


\bibliography{references.bib}

\end{document}
